\chapter{Authentication Mechanics}
\label{ch:authentication}

RGUKT Connect uses stateless token-based authentication. This chapter details password hashing, JSON Web Token (JWT) signatures, request-level verification filters, and WebSocket handshake authentication.

\section{Password Hashing via BCrypt}
User passwords are encrypted before storage. The system uses Spring Security's `BCryptPasswordEncoder` with a default strength factor of 10. BCrypt incorporates a random salt for each hash, securing stored passwords against rainbow table attacks:
\begin{lstlisting}[language=java]
String securePassword = passwordEncoder.encode(requestDTO.getPassword());
user.setPassword(securePassword);
\end{lstlisting}

During login, the system matches the raw password against the stored hash using:
\begin{lstlisting}[language=java]
passwordEncoder.matches(loginRequest.getPassword(), user.getPassword())
\end{lstlisting}

\section{JWT Signature \& Payload Claims}
Upon successful authentication, the server generates a signed JWT access token. 
\begin{itemize}
    \item \textbf{Signature Algorithm:} HMAC-SHA256 (`HS256`).
    \item \textbf{Signing Key:} Loaded from the `JWT\_SECRET` environment variable.
    \item \textbf{Payload Claims:} The JWT contains:
    \begin{itemize}
        \item Subject (`sub`): University email address.
        \item Issued At (`iat`): Timestamp of creation.
        \item Expiration (`exp`): Timestamp when the token expires.
        \item Custom Claims: User ID number and account Role.
    \end{itemize}
\end{itemize}

\section{Request Filtering: `JwtFilter`}
For HTTP requests, the custom `JwtFilter` class (which extends `OncePerRequestFilter`) processes the security context:
\begin{enumerate}
    \item It extracts the `Authorization` header from the incoming request.
    \item It checks if the header begins with the prefix `Bearer `.
    \item It parses the token using the secret key and extracts the email and user roles.
    \item If the token is valid, it builds a `UsernamePasswordAuthenticationToken` and sets it in Spring Security's context:
    \begin{lstlisting}[language=java]
    SecurityContextHolder.getContext().setAuthentication(authenticationToken);
    \end{lstlisting}
\end{enumerate}

\section{WebSocket Handshake Authentication Interceptor}
Since WebSockets are upgrade connections that do not use standard HTTP filters after the initial handshake, securing STOMP frames requires a different approach.

The application registers a custom channel interceptor in `WebSocketConfig` that intercepts incoming STOMP command headers during the connection handshake (`CONNECT` frame):
\begin{enumerate}
    \item The client includes the JWT token in the `Authorization` header inside the STOMP connection headers.
    \item The `configureClientInboundChannel` interceptor extracts the `Authorization` header.
    \item It parses and validates the token.
    \item If valid, it extracts the user's details and registers a custom `Principal` object linked to the user's unique Database ID (converted to a string).
    \item This mapping allows the backend to route private messages to specific users using:
    \begin{lstlisting}[language=java]
    messagingTemplate.convertAndSendToUser(recipientId.toString(), "/queue/messages", messageDto);
    \end{lstlisting}
\end{enumerate}
